\documentclass[Screen16to9,17pt]{foils}
\usepackage{zencurity-slides}

\externaldocument{packet-construction-workshop-exercises}
\selectlanguage{english}
\usepackage{listings}

% Workshop: Getting Started with Network Packets – from ping to 100G

% Do you want to understand what's actually happening down in the layers beneath HTTP? This beginner-friendly workshop takes you from familiar command-line tools like ping and Nmap, through packet manipulation with Python/Scapy, and on to high-performance packet generation with LuaJIT and DPDK via libmoon/MoonGen.

% After a short introductory presentation (~25 slides), you'll get straight to work. Exercises are hands-on and progressive – starting easy and getting more involved as you go. Early exercises include running Nmap to discover hosts on a network and observing ping and ICMP traffic live in Wireshark. From there we build up to constructing and sending your own packets, understanding why this doesn't require low-level C programming, and exploring what it takes to reach wirespeed at 10G – and potentially 100G.

% No prior experience with network programming is required. If you're comfortable with a Linux terminal and curious about what's happening on the wire, you're ready for this workshop.

% Participants should bring: A laptop with a virtual machine running Debian or Kali Linux, with Wireshark, Tcpdump, Python and Nmap installed.

\begin{document}


\mytitlepage{Network Packet Construction Workshop for Beginners}{BornHack July 2026}


\hlkprofiluk

\slide{Goals for today}

\hlkimage{6cm}{bornhack-camp-2024.jpg}

Getting Started with Network Packets -- from Ping to 100G

\begin{list2}
\item Introduce a little basic TCP/IP terminology
\item Let you get some hands on with IP protocols
\item Run ordinary programs like ping, wireshark, tcpdump, Nmap
\item Try more advanced tools like Python with Scapy and Lua
\end{list2}

Photo is NWWC camp at BornHack 2024 \link{https://bornhack.dk}


\slide{Time schedule}f

\hlkimage{4cm}{kame-noanime-small.png}
\begin{list2}
\item 17:00 - 17:45 Introduction and basics
\item 17:45 - 19:00 Play with tools, try sending packets, sniffing packets etc.
\item I will once in a while interrupt with some stuff -- respond to questions for the whole group when something comes up
\item End of workshop I will try to explain where I am with MoonGen a nice tool that was updated in 2025
\end{list2}

\slide{Exercises are completely optional}

We will use a combination of your systems, my networking hardware and my systems.\\
{\color{red} \bf There might be live sniffing done on traffic!}

We try to mimic what you would do in your own networks during the exercises.

Linux is a toolbox I will use and participants are recommended to research virtual machines

Note: even though I talk a lot about Unix and Linux, you can definitely run a lot of tools on Windows and Mac OS X. Today we mostly use command line tools, but some tools have graphical user interfaces.


\slide{Book: Practical Packet Analysis (PPA)}


\hlkimage{6cm}{PracticalPacketAnalysis3E_cover.png}

\emph{Practical Packet Analysis,
Using Wireshark to Solve Real-World Network Problems}
by Chris Sanders, 3rd Edition
April 2017, 368 pp.
ISBN-13:
978-1-59327-802-1
\link{https://nostarch.com/packetanalysis3}

I recommend this book for people new to networking, it has been in HumbleBundle book bundles multiple times


\slide{Prerequisites}

If you are interested in TCP/IP you are welcome

If you want to be an expert in IP and network security I recommend doing exercises

\begin{list1}
\item It is recommended to use virtual machines for the exercises
\item Network security and most internet related security work has the following requirements:
\begin{list2}
\item Network experience
\item TCP/IP principles - often in more detail than a common user
\item Programming is an advantage, for automating things
\item Some Linux and Unix knowledge is in my opinion a {\bf necessary skill} for infosec work\\
-- too many new tools to ignore, and lots found at sites like Github and Open Source written for Linux
\end{list2}
\end{list1}



\slide{Course Network: The BornHack Network}
.
\hlkrightpic{65mm}{-1cm}{sample-network.png}

\begin{list1}
\item I have a course network with me which has the following information:
\begin{list2}
\item Routers and Switches Juniper
\item wireless access-points Aruba
\item IPv4 addresses: 151.216.96.0/20
\item IPv6 addresses: 2001:678:9ec::/48
\item Autonomous system number: AS208647 BornHack\\
\url{https://en.wikipedia.org/wiki/Autonomous_system_(Internet)}

\end{list2}
\end{list1}

\vspace{10mm}
\centerline{\Large You are encouraged to use the network}



\slide{Internet Today}

\hlkimage{10cm}{images/server-client.pdf}

\begin{list1}
\item Clients and servers, roots in the academic world
\item Protocols are old, some more than 30 years
\item Very few protocols where encrypted, today a lot has switched to HTTPS and TLS
\end{list1}


\slide{Network Protocol Knowledge Needed for Network Security}

To work with network security the following protocols are the bare minimum to know about.

\begin{list2}
\item ARP Address Resolution Protocol for IPv4
\item NDP Neighbor Discovery Protocol for IPv6
\item IPv4 \& IPv6 -- the basic packet fields source, destination,
\item ICMPv4 \& ICMPv6 Internet Control Message Protocol
\item UDP User Datagram Protocol
\item TCP Transmission Control Protocol
\item DHCP Dynamic Host Configuration Protocol
\item DNS Domain Name System
\end{list2}

These protocols are part of the Internet Protocol suite, or TCP/IP for short. The canonical document describing this is from 1981 RFC-0791 \citetitle{RFC0791}. The protocols were deployed on the internet around 1983.


\slide{Common Address Space}presentations/network/creative-packets/creative-packets

\vskip 2 cm
\hlkimage{13cm}{IP-address.pdf}

\begin{list2}
\item Internet is defined by the address space
\item IPv4 based on 32-bit addresses, example dotted decimal format 10.0.0.1
\item IPv6 very similar to IPv4 without NAT, 128-bit addresses in hex ::1, 2a06:d380:0:101::80
\end{list2}


\slide{How to use the Internet Protocols (IP)}

Names are used by humans
\begin{center}\hlkbig
www.kramse.org

hlk@kramse.org
\end{center}

Computers use the addresses

\begin{alltt}
www     IN      A       185.129.63.130
        IN      AAAA    2a06:d380:0:102::80
mail    IN      A       217.157.63.250
        IN      AAAA    2a06:d380:0:102::25
\end{alltt}


\slide{Internet ABC}

\begin{list1}
\item Previously we used classes: A, B, C, D og E -- we use CIDR now!
\item This proved to be a bit inflexible:
\begin{list2}
\item A-class has 16 million hosts
\item B-class about 65.000 hosts
\item C-class only 250 hosts
\end{list2}
\item Most people asked for B-class - starting to run out!
\item D-klasse used for multicast
\item E-klasse reserved
\item See \link{http://en.wikipedia.org/wiki/Classful\_network}
\end{list1}

\vskip 5mm
\centerline{\bf Stop saying C, say /24}


\slide{RFC-1918 Private Networks}

\begin{list1}
\item We have some IP addresses anyone can use for private networks
\begin{list2}
\item 10.0.0.0    -  10.255.255.255  (10/8 prefix)
\item 172.16.0.0  -  172.31.255.255  (172.16/12 prefix)
\item 192.168.0.0 -  192.168.255.255 (192.168/16 prefix)
\end{list2}
\item Address Allocation for Private Internets {\bf RFC-1918}
\end{list1}

\begin{alltt}
The blocks 192.0.2.0/24 (TEST-NET-1), 198.51.100.0/24 (TEST-NET-2),
and 203.0.113.0/24 (TEST-NET-3) are provided for use in
documentation.

169.254.0.0/16 has been ear-marked as the IP range to use for end node
auto-configuration when a DHCP server may not be found
\end{alltt}


\slide{Documentation Prefix, IPv6 updates etc.}

Even documentation has its own prefix, RFC5737:
\begin{alltt}
The blocks 192.0.2.0/24 (TEST-NET-1), 198.51.100.0/24 (TEST-NET-2),
and 203.0.113.0/24 (TEST-NET-3) are provided for use in
documentation.
\end{alltt}

IPv6 listed in RFC3849 \verb+2001:DB8::/32+

See RFC3330 \emph{Special-Use IPv4 Addresses} which is updated by
RFC6890 \emph{Special-Purpose IP Address Registries} which in turn is updated by RFC8190\\
Use the web version of RFCs to surf back and forth \link{https://www.rfc-editor.org/rfc/rfc8190}


\slide{CIDR Classless Inter-Domain Routing}

\hlkimage{15cm}{CIDR-aggregation.pdf}

\begin{list2}
\item Subnet mask originally inferred by the class
\item Started to allocate multiple C-class networks - save remaining B-class\\
Resulted in routing table explosion - btw Stop using A, B, C
\item A subnet mask today is a row of 1-bit
\item Supernet, supernetting
\item 10.0.0.0/24 means the network 10.0.0.0 with 24 subnet bits (mask 255.255.255.0)
\item 2a06:d380:0:101::80/64 means the network with 64-bit prefix length
\end{list2}



\slide{Protocols: OSI and Internet models}

\hlkimage{11cm,angle=90}{images/compare-osi-ip.pdf}


\slide{Ethernet, cables}

\hlkimage{17cm}{ethernetLights.jpg}

\centerline{Show link, and activity -- blinkenlights}


\slide{A modern router}

\hlkimage{7cm}{sft1200_1.png}

\begin{list2}
\item Opal (GL-SFT1200) is a pocket-sized travel router supporting 1200Mbps wireless,
Max. 300 Mbps (2.4GHz) + 867 Mbps (5GHz) Fast Wi-Fi Speeds, Powerful CPU with 3 x Gigabit Ports,
Excellent Security with VPN -- OpenVPN \& WireGuard, IPv6
\end{list2}
Source: \url{https://www.gl-inet.com/products/gl-sft1200/}


\slide{Packets across the wire or wireless}

\hlkimage{20cm}{ethernet-frame-1.pdf}
\begin{list1}
\item Looking at data as a stream the packets are a pattern laid on top
\item Network technology defines the start and end of a frame, example Ethernet
\item From a lower level we receive a packet, example 1500-bytes from Ethernet driver
\item Operating system masks a lot of complexity
\end{list1}


\slide{The basic tools today}

{\Large Knowledge and insight}
\begin{list2}
\item Networks have end-points and conversations on multiple layers
\item Wireshark is advanced, try right-clicking different places
\item Name resolution includes low level MAC addresses, and IP - names
\end{list2}

\begin{list2}
\item Tcpdump format, built-in to many network devices
\item Remote packet dumps, like \verb+tcpdump –i eth0 –w packets.pcap+
\item Story: tcpdump was originally written in 1988 by Van Jacobson, Sally Floyd, Vern Paxson and Steven McCanne who were, at the time, working in the Lawrence Berkeley Laboratory Network Research Group\\
 \link{https://en.wikipedia.org/wiki/Tcpdump}
\end{list2}

\vskip 5mm

\centerline{\Large Great network security comes from knowing networks!}


\slide{IPv4 header - RFC-791 september 1981}

\begin{alltt}\footnotesize
    0                   1                   2                   3
    0 1 2 3 4 5 6 7 8 9 0 1 2 3 4 5 6 7 8 9 0 1 2 3 4 5 6 7 8 9 0 1
   +-+-+-+-+-+-+-+-+-+-+-+-+-+-+-+-+-+-+-+-+-+-+-+-+-+-+-+-+-+-+-+-+
   |Version|  IHL  |Type of Service|          Total Length         |
   +-+-+-+-+-+-+-+-+-+-+-+-+-+-+-+-+-+-+-+-+-+-+-+-+-+-+-+-+-+-+-+-+
   |         Identification        |Flags|      Fragment Offset    |
   +-+-+-+-+-+-+-+-+-+-+-+-+-+-+-+-+-+-+-+-+-+-+-+-+-+-+-+-+-+-+-+-+
   |  Time to Live |    Protocol   |         Header Checksum       |
   +-+-+-+-+-+-+-+-+-+-+-+-+-+-+-+-+-+-+-+-+-+-+-+-+-+-+-+-+-+-+-+-+
   |                       Source Address                          |
   +-+-+-+-+-+-+-+-+-+-+-+-+-+-+-+-+-+-+-+-+-+-+-+-+-+-+-+-+-+-+-+-+
   |                    Destination Address                        |
   +-+-+-+-+-+-+-+-+-+-+-+-+-+-+-+-+-+-+-+-+-+-+-+-+-+-+-+-+-+-+-+-+
   |                    Options                    |    Padding    |
   +-+-+-+-+-+-+-+-+-+-+-+-+-+-+-+-+-+-+-+-+-+-+-+-+-+-+-+-+-+-+-+-+

                    Example Internet Datagram Header
\end{alltt}
Source: \url{https://datatracker.ietf.org/doc/html/rfc791} and updated later

\slide{IPv6 header - RFC-1883 December 1995}
\vskip -1cm
\begin{alltt}\footnotesize
   +-+-+-+-+-+-+-+-+-+-+-+-+-+-+-+-+-+-+-+-+-+-+-+-+-+-+-+-+-+-+-+-+
   |Version| Traffic Class |           Flow Label                  |
   +-+-+-+-+-+-+-+-+-+-+-+-+-+-+-+-+-+-+-+-+-+-+-+-+-+-+-+-+-+-+-+-+
   |         Payload Length        |  Next Header  |   Hop Limit   |
   +-+-+-+-+-+-+-+-+-+-+-+-+-+-+-+-+-+-+-+-+-+-+-+-+-+-+-+-+-+-+-+-+
   |                                                               |
   +                                                               +
   |                                                               |
   +                         Source Address                        +
   |                                                               |
   +                                                               +
   |                                                               |
   +-+-+-+-+-+-+-+-+-+-+-+-+-+-+-+-+-+-+-+-+-+-+-+-+-+-+-+-+-+-+-+-+
   |                                                               |
   +                                                               +
   |                                                               |
   +                      Destination Address                      +
   |                                                               |
   +                                                               +
   |                                                               |
   +-+-+-+-+-+-+-+-+-+-+-+-+-+-+-+-+-+-+-+-+-+-+-+-+-+-+-+-+-+-+-+-+
\end{alltt}
%Source: \url{https://datatracker.ietf.org/doc/html/rfc1883} and obsoleted by\\
%\url{https://datatracker.ietf.org/doc/html/rfc8200}

\slide{UDP User Datagram Protocol}
\hlkimage{16cm}{udp-1.pdf}
Connectionless \url{https://en.wikipedia.org/wiki/User_Datagram_Protocol}


\slide{TCP Transmission Control Protocol}
\hlkimage{14cm}{tcp-1.pdf}

Connection-oriented \url{https://en.wikipedia.org/wiki/Transmission_Control_Protocol}


\slide{TCP three way handshake}

\hlkimage{6cm}{images/tcp-three-way.pdf}

\begin{list2}
\item Session setup is used in some protocols
\item Other protocols like HTTP/2 can perform request in the first packet
\end{list2}


\slide{Well-Known Port Numbers}
presentations/network/creative-packets/creative-packets
\hlkimage{3cm}{iana1.jpg}

\begin{list1}
\item IANA maintains a list of magical numbers in TCP/IP
\item Lists of protocl numbers, port numers etc.
\item A few notable examples:
\begin{list2}
\item Port 25/tcp Simple Mail Transfer Protocol (SMTP)
\item Port 53/udp and 53/tcp Domain Name System (DNS)
\item Port 80/tcp Hyper Text Transfer Protocol (HTTP)
\item Port 443/tcp HTTP over TLS/SSL (HTTPS)
\end{list2}
\item Source: \link{http://www.iana.org}
\end{list1}


\slide{ICMP and Ping}

\begin{list1}
\item Internet Control Message Protocol (ICMP)
\item Error conditions are signalled using this
\item The ping program sends ICMP ECHO request and
 expect ICMP ECHO reply
\item
\begin{alltt}
\small {\bfseries
$ ping 192.168.1.1}
PING 192.168.1.1 (192.168.1.1): 56 data bytes
64 bytes from 192.168.1.1: icmp_seq=0 ttl=150 time=8.849 ms
64 bytes from 192.168.1.1: icmp_seq=1 ttl=150 time=0.588 ms
64 bytes from 192.168.1.1: icmp_seq=2 ttl=150 time=0.553 ms
\end{alltt}
\end{list1}

This is our first \emph{packet construction tool}

\slide{Ping with options}

Ping builds packets with the ICMP protocol and usually send a basic ICMP ECHO Request, see {https://datatracker.ietf.org/doc/html/rfc792} and updates
\begin{alltt}
Echo or Echo Reply Message

    0                   1                   2                   3
    0 1 2 3 4 5 6 7 8 9 0 1 2 3 4 5 6 7 8 9 0 1 2 3 4 5 6 7 8 9 0 1
   +-+-+-+-+-+-+-+-+-+-+-+-+-+-+-+-+-+-+-+-+-+-+-+-+-+-+-+-+-+-+-+-+
   |     Type      |     Code      |          Checksum             |
   +-+-+-+-+-+-+-+-+-+-+-+-+-+-+-+-+-+-+-+-+-+-+-+-+-+-+-+-+-+-+-+-+
   |           Identifier          |        Sequence Number        |
   +-+-+-+-+-+-+-+-+-+-+-+-+-+-+-+-+-+-+-+-+-+-+-+-+-+-+-+-+-+-+-+-+
   |     Data ...
   +-+-+-+-+-
\end{alltt}

Lets try options \verb+-t+ (ttl) \verb+-s+ and possibly \verb+-f+


\slide{traceroute}

\begin{list1}
  \item traceroute works using the Time to Live (TTL) or Hop-Count fields
\item Each router will subtract one from this, and if zero -- return ICMP message
\item Unix uses UDP and Windows usually uses ICMP Echo
\item
\begin{alltt}
\small{\bfseries
$ traceroute 185.129.60.129}
traceroute to 185.129.60.129 (185.129.60.129),
30 hops max, 40 byte packets
 1  safri (10.0.0.11)  3.577 ms  0.565 ms  0.323 ms
 2  router (185.129.60.129)  1.481 ms  1.374 ms  1.261 ms
\end{alltt}
\end{list1}


\slide{Traceroute}

You know traceroute, the hacker tool?

\begin{enumerate}
\item Create UDP packet with $TTL = 1$
\item Send packet, get response from closest router
\item Repeat three times while increasing port number
\item increase TTL and repeat, until destination reached or max TTL
\end{enumerate}
\vskip 1 cm

\centerline{\Large\bf Not a built-in feature of IP/TCP}
\vskip 1 cm
(Microsoft Windows uses ICMP and calls the tools tracert)



\slide{IPv6 addressing RFC-4291}

\begin{list1}
\item Addresses are always 128-bit identifiers for interfaces and sets of
   interfaces
\item Unicast:   An identifier for a {\bf single interface}.\\
A packet sent to a
               unicast address is delivered to the interface identified
               by that address.
\item Anycast:   An identifier for a {\bf set of interfaces} (typically
               belonging to different nodes).\\  A packet sent to an
               anycast address is {\bf delivered to one} of the interfaces
               identified by that address (the "nearest" one, according
               to the routing protocols' measure of distance).

\item Multicast: An identifier for a {\bf set of interfaces} (typically
               belonging to different nodes). \\ A packet sent to a
               multicast address is {\bf delivered to all interfaces
               identified by that address}.
\end{list1}

\slide{IPv6 addressing RFC-4291, cont.}

\hlkimage{22cm}{ipv6-address-1.pdf}

\begin{list1}
\item 8 times 4 hex-digits seperated by colon x:x:x:x:x:x:x:x
\item Written as ipv6-address/prefix-length CIDR notation
\item Leading zeros can be removed
\item One or more groups of 16 bits of zeros can be replaced by ::
\end{list1}


\slide{Examples:}
\begin{list2}
\item ABCD:EF01:2345:6789:ABCD:EF01:2345:6789

\item Adddress 2001:DB8:0:0:8:800:200C:417A
\item Address of loopback ::1
\item IPv6 prefix 2a02:09d0:95::1/64, subnet 2a02:09d0:0095:0000::/64
\item Address 2a02:09d0:95::1 or 2a02:09d0:0095:0000:0000:0000:0000:0001
\end{list2}



\slide{IPv6 address - special prefixes }

\begin{list2}
\item link-local unicast addresses\\
fe80::/10 generated from the interface MAC address EUI-64
\item FEC0::/10 site-local - deprecated in RFC-3879

\item 2001:0DB8::/32 NON-ROUTABLE range to be used for documentation purpose RFC-3849.

\item FC00::/7 Unique Local IPv6 Unicast Addresses RFC-4193\\
\link{http://www.simpledns.com/private-ipv6.aspx}\\
If you do not like to put public addresses on internal network - use this instead
\end{list2}



\slide{Hello neighbors}

\begin{alltt}\small
$ ping6 -w -I en1 ff02::1
PING6(72=40+8+24 bytes) fe80::223:6cff:fe9a:f52c%en1 --> ff02::1
30 bytes from fe80::223:6cff:fe9a:f52c%en1: bigfoot
36 bytes from fe80::216:cbff:feac:1d9f%en1: mike.kramse.dk.
38 bytes from fe80::200:aaff:feab:9f06%en1: xrx0000aaab9f06
34 bytes from fe80::20d:93ff:fe4d:55fe%en1: harry.local
36 bytes from fe80::200:24ff:fec8:b24c%en1: kris.kramse.dk.
31 bytes from fe80::21b:63ff:fef5:38df%en1: airport5
32 bytes from fe80::216:cbff:fec4:403a%en1: main-base
44 bytes from fe80::217:f2ff:fee4:2156%en1: Base Station Koekken
35 bytes from fe80::21e:c2ff:feac:cd17%en1: arnold.local
\end{alltt}



\slide{Wireshark - graphical network sniffer}

\hlkimage{13cm}{images/wireshark-http.png}

\centerline{Capture - Options, select a network interface}
\centerline{\link{http://www.wireshark.org}}


\slide{Your Privacy }

\hlkimage{18cm}{images/internet-browsing.pdf}

\begin{list2}
\item Your data travels far
\item Often crossing borders, virtually and literally
\end{list2}


\slide{Basic port scanning}

\begin{list1}
\item What is a port scan
\item Testing all values possible for port number from 0/1 to 65535
\item Goal is to identify open ports, listening and vulnerable services
\item Most often TCP og UDP scan
\item TCP scanning is more realiable than UDP scanning
\item TCP handshake must respond with SYN-ACK packets
\item UDP applications respond differently -- if they even respond\\
so probes with real requests may get response, no firewall they respond withb ICMP on closed ports
\item Use the GUI program Zenmap while learning Nmap
\end{list1}

\slide{TCP three-way handshake}

\hlkimage{4cm}{images/tcp-three-way.pdf}

\begin{list2}
\item {\bfseries TCP SYN half-open} scans
\item in the old days systems would only log when a full TCP connection was setup\\
  -- so doing only half open it was a \emph{stealth}-scans
\item Today system and IDS intrusion detection can easily monitor for this
\item Sending a lot of SYN packets can create a Denial of Service -- {\bfseries SYN-flooding}
\end{list2}
\slide{Scope: select systems for testing}

\hlkimage{10cm}{overview-routing-customer-2015.png}

\begin{list2}
\item Routers in front of critical systems and networks - availability
\item Firewalls -- are traffic flows restricted
\item Mail servers -- open for relaying
\item Web servers -- remote code execution in web systems, data download
\end{list2}

\slide{Ping and port sweep}

\begin{list1}
\item Scans across the network are named sweeps
\item Ping sweeps using ICMP Ping probes
\item Port sweep trying to find a specific service, like port 80 web
\item Quite easy to see in network traffic:
\begin{list2}
\item Selecting two IP-adresser not in use
\item Should not see any traffic, but if it does, its being scanned
\item If traffic is received on both addresses, its a sweep -- if they are a bit apart it is even better, like 10.0.0.100 and 10.0.0.200
  \end{list2}

\vskip 2cm
Pro tip: a Great network intrusion detection engine (IDS), is Suricata \link{suricata-ids.org}
\end{list1}

\slide{what is Nmap today}
\begin{quote}
Nmap ("Network Mapper"presentations/network/creative-packets/creative-packets) is a free and open source (license) utility for network discovery and security auditing.
\end{quote}

\begin{list1}
\item Initial release September 1997; almost 30 years ago
\item Today a package of programs for Windows, Mac, BSD, Linux, ... source
\item Flexible, powerful, and free! Includes other tools!
\end{list1}

Bonus info: you can help Nmap by submitting fingerprints


\slide{Nmap port sweep for web servers}

\begin{alltt}\small
root@cornerstone:~#{\bfseries  nmap -p80,443 172.29.0.0/24}

Starting Nmap 6.47 ( http://nmap.org ) at 2015-02-05 07:31 CET
Nmap scan report for 172.29.0.1
Host is up (0.00016s latency).
PORT    STATE    SERVICE
{\color{darkgreen}80/tcp  open     http}
443/tcp filtered https
MAC Address: 00:50:56:C0:00:08 (VMware)

Nmap scan report for 172.29.0.138
Host is up (0.00012s latency).
PORT    STATE  SERVICE
{\color{darkgreen}80/tcp  open   http}
443/tcp closed https
MAC Address: 00:0C:29:46:22:FB (VMware)

\end{alltt}

\slide{Nmap port sweep for SNMP port 161/UDP}

\begin{alltt}\small
root@cornerstone:~#{\bfseries nmap -sU -p 161 172.29.0.0/24}
Starting Nmap 6.47 ( http://nmap.org ) at 2015-02-05 07:30 CET
Nmap scan report for 172.29.0.1
Host is up (0.00015s latency).
PORT    STATE         SERVICE
{\color{darkgreen}161/udp open|filtered snmp}
MAC Address: 00:50:56:C0:00:08 (VMware)

Nmap scan report for 172.29.0.138
Host is up (0.00011s latency).
PORT    STATE  SERVICE
{\bf{161/udp closed snmp}}
MAC Address: 00:0C:29:46:22:FB (VMware)
...

Nmap done: 256 IP addresses (5 hosts up) scanned in 2.18 seconds
\end{alltt}

\vskip 5mm
\centerline{More reliable to use Nmap script with probes like --script=snmp-info}

\slide{Nmap Advanced OS detection}
\begin{alltt}\footnotesize
root@cornerstone:~#{\bfseries nmap -A -p80,443 172.29.0.0/24}
Starting Nmap 6.47 ( http://nmap.org ) at 2015-02-05 07:37 CET
Nmap scan report for 172.29.0.1
Host is up (0.00027s latency).
PORT    STATE    SERVICE VERSION
80/tcp  open     http    Apache httpd 2.2.26 ((Unix) DAV/2 mod_ssl/2.2.26 OpenSSL/0.9.8zc)
|_http-title: Site doesn't have a title (text/html).
443/tcp filtered https
MAC Address: 00:50:56:C0:00:08 (VMware)
Device type: media device|general purpose|phone
Running: Apple iOS 6.X|4.X|5.X, Apple Mac OS X 10.7.X|10.9.X|10.8.X
OS details: Apple iOS 6.1.3, Apple Mac OS X 10.7.0 (Lion) - 10.9.2 (Mavericks)
or iOS 4.1 - 7.1 (Darwin 10.0.0 - 14.0.0), Apple Mac OS X 10.8 - 10.8.3 (Mountain Lion)
or iOS 5.1.1 - 6.1.5 (Darwin 12.0.0 - 13.0.0)
OS and Service detection performed.
Please report any incorrect results at http://nmap.org/submit/
\end{alltt}

\begin{list2}
\item Low-level way to identify operating systems, also try/use
  \verb+nmap -A+
\item Send probes and observe responses, lookup in table of known OS and responses
\item Techniques known since at least: \emph{ICMP Usage In Scanning} Version 3.0,
  Ofir Arkin, 2001 %\link{https://web.archive.org/web/20050210093427/http://www.sys-security.com/html/projects/icmp.html} % Original side er død
\end{list2}


\slide{Juniper DoS}

\begin{quote}
	JUNOS (Juniper) Flaw Exposes Core Routers to Kernel Crash\\
A report has been received from Juniper at 4:25pm under bulletin PSN-2010-01-623 that {\bf a crafted malformed TCP field option in the TCP header of a packet will cause the JUNOS kernel to core (crash).} In other words the kernel on the network device (gateway router) will crash and reboot if a packet containing this crafted option is received on a listening TCP port. The JUNOS firewall filter is unable to filter a TCP packet with this issue. Juniper claims this issue as exploit was identified during investigation of a vendor interoperability issue.
\end{quote}

Juniper bulletin PSN-2010-01-623 and \\
{\footnotesize http://praetorianprefect.com/archives/2010/01/
junos-juniper-flaw-exposes-core-routers-to-kernal-crash/}



\slide{History - fast forward}

\begin{list1}
\item Fast forward, another day or over beer
\item Various snoop and sniffer programs
\item Libnet - generic API to help with the construction and handling of network packets
\item Lidpcap
\end{list1}

\centerline{tl;dr - tools became portable or implemented on Linux}

\vskip 3cm
Hint: if not already done, create a virtual machine with\\
 \link{http://www.backtrack-linux.org/}


\slide{icmpush, Libnet}

\begin{list1}
\item icmpush - ICMP packets galore
\item Nemesis can natively craft and inject ARP, DNS, ETHERNET, ICMP, IGMP, IP, OSPF, RIP, TCP and UDP packets
\link{http://nemesis.sourceforge.net}

% libnet not working, empty \link{http://code.google.com/p/libnet/}
Archive of packetfactory exist at \link{http://packetfactory.openwall.net/}

\item Paketto Keiretsu author Dan Kaminsky, circa 2002
\end{list1}


\slide{Command Line packet generation}

\begin{list1}
\item hping is a command-line oriented TCP/IP packet assembler/analyzer\\ \link{http://www.hping.org/}
\item  ISIC is a suite of utilities to exercise the stability of an IP Stack and its component stacks (TCP, UDP, ICMP et. al.) \link{http://isic.sourceforge.net/}
\item TCP traceroute programs - make sure you try out the Nping from the Nmap suite
\item Too many tools to mention
\end{list1}

\slide{SING stands for 'Send ICMP Nasty Garbage'}

SING stands for 'Send ICMP Nasty Garbage'. It is a tool that sends ICMP packets fully customized from command line. Its main purpose is to replace the ping command but adding certain enhancements (Fragmentation, spoofing,...)

\link{http://sourceforge.net/projects/sing/}

\vskip 2 cm
\centerline{\bf\Large Please do good with the tools :-)}
\vskip 2 cm
Maybe we should do some network wars with packets :-)

\slide{Perl and Python libraries}

\begin{list1}
\item You can program packets in C, but C respects you enough to let you shoot yourself in the foot
\item Perl has a multitude of nice libraries, some create packets - Net:: Packet and Net::DNS
\item Today I would recommend Python with Scapy and Nmap with Lua for beginners
\item The Go Language also has some libraries for this
\item Choose your own adventure and languages
\end{list1}


\slide{Exploit code / proof of concept}

\inputminted[fontsize=\small,  breaklines]{perl}{files/junos-crash.pl}
{\small\link{http://evilrouters.net/2010/01/09/junos-psn-2010-01-623-exploit/}}


\slide{Use scapy to send JunOS killin' packet}

%\hlkimage{16cm}{files/scapy-junos-crash.png}
\inputminted[fontsize=\small,  breaklines]{python}{files/junos-crash.py}
{\small
\link{http://evilrouters.net/2010/01/10/use-scapy-to-send-junos-killin-packet/}}\\
+ comment about TCP options


\slide{Scapy}
% my tools /userdata/network/scapy

\begin{alltt}
>>> for i in range (5,10):
...       send(IP(dst="192.168.0.1")/ICMP())
...
.
Sent 1 packets.
.
Sent 1 packets.
.
Sent 1 packets.
.
Sent 1 packets.
.
Sent 1 packets.
\end{alltt}

\slide{Save the packets}
It is often useful to save capture packets to pcap file for use at later time or with different applications:

\begin{alltt}
>>> wrpcap("temp.cap",pkts)
To restore previously saved pcap file:

>>> pkts = rdpcap("temp.cap")
or

>>> pkts = sniff(offline="temp.cap")
\end{alltt}

This is basic functionality with Scapy

Maybe start in the Scapy documentation and run some examples there!

\slide{Speeding up}

Scapy will happily build your packet. It just won't build it fast.
\begin{list1}
\item Why generate and send, when you can send faster by pre-generating:
\item  tcpreplay \link{http://tcpreplay.synfin.net/}
\item  pcapdiff \link{http://www.eff.org/testyourisp/pcapdiff/}
\item  main observation, there are already lots of tools
\end{list1}

\slide{More speed}

\begin{list1}
\item Zero copy, stop moving those bits!
\item Performance doc file from \link{http://netsniff-ng.org/}
\end{list1}

\begin{alltt}\small
[1] http://www.linuxfoundation.org/collaborate/workgroups/networking/napi
[2] http://datatag.web.cern.ch/datatag/howto/tcp.html
[3] http://thread.gmane.org/gmane.linux.network/191115
    Kernel build option:
      CONFIG_HAVE_BPF_JIT=y
      CONFIG_BPF_JIT=y
[4] http://bit.ly/3XbBrM
\end{alltt}

A lot of things have happened over the years!

\slide{netsniff-ng}

The netsniff-ng toolkit consists of the following utilities:

\begin{list2}
\item netsniff-ng, a zero-copy analyzer, pcap capturer and replayer
\item trafgen, a high-performance zero-copy network traffic generator
\item bpfc, a Berkeley Packet Filter compiler supporting Linux extensions
\item ifpps, a top-like kernel networking and system statistics tool
\item flowtop, a top-like netfilter connection tracking tool
\item curvetun, a lightweight multiuser IP tunnel based on elliptic curve cryptography
\item ashunt, an Autonomous System (AS) trace route and ISP testing utility
\end{list2}


http://netsniff-ng.org/


\slide{Programmable cards}

\hlkimage{10cm}{files/NetFPGA10G_front_b.jpg}

\begin{list1}
\item future programmable cards  - high speed possible, but more complex programming
\item \link{http://netfpga.org/foswiki/bin/view/NetFPGA/WebHome}
\item \link{https://github.com/crotsos/netfpga-packet-generator-c-library}
\item NetFPGA-10G shown in picture \emph{The Academic price is \$1,675.}
\end{list1}

I dreamed of this card, but did NOT have the money for it.


\slide{Penguinping testing}

\hlkimage{16cm}{ddos-test-layout.png}

\slide{Penguinping Software stack}

\begin{itemize}
\item I used a nice library Libmoon and MoonGen for creating packets with Lua and 10G cards -- one CPU core could push 14million packets per second
\item Original paper: P.~Emmerich et al., \emph{MoonGen: A Scriptable High-Speed Packet Generator}, ACM IMC 2015. \url{https://dl.acm.org/doi/10.1145/2815675.2815692}
\item Updated fork with MoonEm path emulator (delay, rate limiting, packet loss): TU Munich Network Architectures and Services group, CoNEXT 2025. \url{https://github.com/tumi8/moongen}
\item libmoon (core DPDK/LuaJIT framework underlying MoonGen): \url{https://github.com/libmoon/libmoon}
\item This is confirmed working on Debian 13 in 2026!
\item It used DPDK, Just-in-time compilation and is very nice
\end{itemize}

I promise to update \url{https://penguinping.org/} and \url{https://codeberg.org/kramse/penguinping}

\slide{Nvidia Mellanox DPU2}

\hlkimage{12cm}{nvidia-dpu2.jpg}

\begin{itemize}
\item MT42822 — ConnectX-6 Dx network controller PCIe Gen4 host interface
\item 2× 100G ports (SFP28/DAC)
\item 8× Arm Cortex-A72 onboard cores
\item Drivers: \verb+mlx5_core+ (host), \verb+mlx5+ PMD (DPDK)
 \end{itemize}

\slide{Packet Generation — MoonGen on BlueField-2}
\begin{itemize}
\item Host: Dell PowerEdge T150, Intel Xeon E-2314 (4 cores, 2.80 GHz)
\item NIC: NVIDIA BlueField-2, single 100G port via DAC loopback
\item Software: MoonGen + libmoon, DPDK \verb+mlx5+ PMD
\item Packet size: 60 bytes (minimum Ethernet frame)
\item TX queues: 3 (1 core reserved for DPDK master)
\item \textbf{Result: \textasciitilde{}43 Mpps / \textasciitilde{}20 Gbps}
\item Verified via \verb+ethtool -S+ hardware counters (\verb+rx_packets_phy+)
\item Bottleneck: CPU cores — scales linearly with queue count
\end{itemize}

It's in my tent on PTY, if you want to try it! aaaand I brought a 10-core CPU PC

\slide{Conclusion -- Networks are exciting!}

\hlkimage{10cm}{IMG_20250309_bailey.jpg}


\begin{list2}
\item Many protocols work together to provide \emph{cat and dog pictures}
\item Lots of acronyms, but a relatively small set of protocols will get you started
\item Working with a small portable router and your laptop might teach you the basics
\end{list2}

%Pictured a small white and black dog playing with a toy porcupine
\myquestionspage



\slide{Advanced Subjects}

Below are some pointers to very advanced tools.

They are listed here as a next step from doing manual tcpdump and Wireshark packet capture.

The tools also allow you to continuously monitor a network easily.

\begin{list2}
\item Zeek
\item Suricata
\end{list2}

That said, they are also great learning tools -- start from a packet capture files, decode it, get an overview and read about it.

\slide{Architecture for packet capture}

\hlkimage{4cm}{network-horiz-onion.png}
Source: picture from \link{https://docs.securityonion.net/en/2.3/introduction.html}

\begin{list2}
\item Note the terminology North-South -- from the internet into the systems
\item East-West -- horizontal traffic inside the data center
\item See also from Security Onion \url{https://docs.securityonion.net/en/2.3/architecture.html#architecture}
\end{list2}



\slide{The Zeek Network Security Monitor}

\hlkimage{13cm}{zeek-overview.png}

The Zeek Network Security Monitor is not a single tool, more of a
powerful network analysis framework. Note: the project was renamed from Bro to Zeek in Oct 2018

{\small Source \url{https://www.zeek.org/}}


\slide{Suricata IDS/IPS/NSM}

\hlkimage{6cm}{suricata.png}

\begin{quote}
Suricata is a high performance Network IDS, IPS and Network Security Monitoring engine. Open Source and owned by a community run non-profit foundation, the Open Information Security Foundation (OISF). Suricata is developed by the OISF and its supporting vendors.
\end{quote}

 \link{http://suricata-ids.org/}
 \link{http://openinfosecfoundation.org}

\centerline{I often use Suricata and Zeek together}


\end{document}
